\section{Limitations and Responsible Use}
Private validation partitions can differ systematically across clients, so median
quality is not an unbiased estimator of global utility. Fixed-panel monitoring
reduces one source of score variation but cannot distinguish every benign drift
from poisoning. Detector error rates and utility before and after recovery must be
interpreted together.

Checkpoint restoration can lose useful progress and leave earlier poison intact.
Wrongful penalties can exclude honest validators and reduce future committee
availability. Stake requirements also create participation barriers. These effects
limit any claim that staking alone makes an attack unprofitable. A utility model
and strategic-adversary analysis would be needed for such a claim.

The local EVM provides contract execution and transaction receipts on one host.
It does not measure network partitions, distributed consensus, adversarial block
production, or production checkpoint availability. The trusted registrar and
orchestrator are material deployment assumptions. Metric releases and individual
model updates retain privacy risks, and an anonymous deployment would require a
separate enrollment mechanism.

Classification findings require separate validation for search routing.
The new SIFT1M labels have explicit provenance, but the case-study outcomes
remain pending. Their sharding differs from the original pilot, so numerical
comparisons with its reported recall would be invalid.
